% Vertical timeline for two-column DnD guide layout.
%
% Usage:
%   \begin{DndTimeline}[The Fade at a Glance]
%     \DndTimelineEvent{3850 SC}{First stirrings of the death of magic.}
%     \DndTimelineEvent{4218 SC}{Present year.}
%   \end{DndTimeline}

\RequirePackage{tikz}
\usetikzlibrary{calc}

\newcounter{dndtimelineevent}

% TikZ x units = cm. Spine at x=0; dates/text at this offset.
\newcommand{\dndtimelinecontent}{0.14}

\NewDocumentEnvironment{DndTimeline}{O{}}{%
  \setcounter{dndtimelineevent}{0}%
  \IfValueT{#1}{%
    \par\vspace{0.25\baselineskip}%
    \begingroup\DndFontTableTitle #1\nopagebreak\par\endgroup
    \vspace{0.35\baselineskip}%
  }{}%
  \begin{center}%
  \begin{tikzpicture}[
    x=1cm,
    y=1cm,
    dnd timeline dot/.style={fill=titlered, circle, inner sep=1.75pt},
    dnd timeline date/.style={anchor=west, font=\bfseries\small},
    dnd timeline text/.style={
      anchor=north west,
      align=left,
      text width=\dimexpr\columnwidth-0.75cm\relax,
      font=\small,
    },
  ]
}{%
  \ifnum\value{dndtimelineevent}>0\relax
    \draw[titlered, line width=0.9pt]
      (dndtimelinefirstdot) -- (dndtimelinelastdot);
  \fi
  \end{tikzpicture}%
  \end{center}%
  \vspace{0.35\baselineskip}%
}

\newcommand{\DndTimelineEvent}[2]{%
  \stepcounter{dndtimelineevent}%
  \edef\dndtimelineid{\number\value{dndtimelineevent}}%
  \def\dndtimelinedate{#1}%
  \def\dndtimelinetext{#2}%
  \ifnum\value{dndtimelineevent}=1\relax
    \node[dnd timeline date] (dndt-\dndtimelineid-d)
      at (\dndtimelinecontent, 0) {\dndtimelinedate};
  \else
    \node[dnd timeline date] (dndt-\dndtimelineid-d)
      at ($ (dndtimelineprev) + (0,-0.55) $) {\dndtimelinedate};
  \fi
  \node[dnd timeline text] (dndt-\dndtimelineid-v)
    at ([yshift=-0.15cm]dndt-\dndtimelineid-d.south west) {\dndtimelinetext};
  \node[dnd timeline dot] (dndt-\dndtimelineid-o)
    at (0, 0 |- dndt-\dndtimelineid-d.west) {};
  \ifnum\value{dndtimelineevent}=1\relax
    \coordinate (dndtimelinefirstdot) at (dndt-\dndtimelineid-o);
  \fi
  \coordinate (dndtimelinelastdot) at (dndt-\dndtimelineid-o);
  \coordinate (dndtimelineprev) at (dndt-\dndtimelineid-v.south west);
}
